\section{Architecture technique}

\subsection{Vue d'ensemble}

Wysp tourne entièrement en local. Un processus Python, \textbf{Wysp Core}, sert le \textbf{panneau} (application React ouverte dans une fenêtre Chrome à côté du logiciel) et accepte les connexions des \textbf{adaptateurs}. Chaque adaptateur vit dans son application hôte : une extension Blender, une extension Chrome. Le cœur appelle un fournisseur LLM et un service de synthèse vocale configurés par variables d'environnement.

\begin{figure}[h]
\centering
\begin{tikzpicture}[
  font=\sffamily\scriptsize,
  box/.style={draw=#1, fill=#1!8, rounded corners=2pt, align=center, minimum height=0.62cm, inner sep=3pt},
  mod/.style={draw=pa!70, fill=white, rounded corners=1.5pt, align=center, minimum width=2.05cm, minimum height=0.55cm, inner sep=2pt},
  lab/.style={font=\sffamily\scriptsize\itshape, text=muted},
  arr/.style={-{Stealth[length=2.2mm]}, thick, color=surface},
  biarr/.style={{Stealth[length=2.2mm]}-{Stealth[length=2.2mm]}, thick, color=surface}]

  % Core
  \node[draw=pa, fill=pa!5, rounded corners=4pt, minimum width=9.0cm, minimum height=4.5cm] (core) at (0,0.5) {};
  \node[font=\sffamily\bfseries\small, text=pa] at (0,2.47) {Wysp Core (Python, FastAPI, asyncio)};
  \node[draw=pa!70, fill=pa!12, rounded corners=1.5pt, minimum width=8.7cm, minimum height=0.42cm, inner sep=1pt] (ctr) at (0,1.98) {contract : modèles Pydantic, prédicats, empreinte d'état (partagés avec les SDK et le kit)};
  \node[mod, minimum width=1.5cm, minimum height=2.85cm] (gw) at (-3.65,0.12) {gateway\\[2pt]\tiny JSON-RPC\\\tiny registre\\\tiny jetons};
  \node[mod, minimum width=1.5cm, minimum height=2.85cm] (ui) at (3.65,0.12) {ui\_api\\[2pt]\tiny /ws/ui\\\tiny /api/tts\\\tiny panneau};
  \node[mod, minimum width=1.7cm] (sess)  at (-1.9, 1.12) {session\\\tiny bus d'événements};
  \node[mod, minimum width=1.7cm] (orch)  at ( 0.0, 1.12) {orchestrator\\\tiny tours, outils};
  \node[mod, minimum width=1.7cm] (board) at ( 1.9, 1.12) {board\\\tiny projection};
  \node[mod, minimum width=1.7cm] (trust) at (-1.9, 0.12) {trust\\\tiny porte, undo, audit};
  \node[mod, minimum width=1.7cm] (verif) at ( 0.0, 0.12) {verification\\\tiny prédicats, preuves};
  \node[mod, minimum width=1.7cm] (narr)  at ( 1.9, 0.12) {narration\\\tiny segments};
  \node[mod, minimum width=1.7cm] (ped)   at (-1.9,-0.88) {pedagogy\\\tiny leçons, apprenant};
  \node[mod, minimum width=1.7cm] (llm)   at ( 0.0,-0.88) {llm\\\tiny fournisseurs};
  \node[mod, minimum width=1.7cm] (store) at ( 1.9,-0.88) {store\\\tiny \textasciitilde/.wysp};

  % Adapters
  \node[box=pc, minimum width=2.6cm, minimum height=1.35cm] (bl) at (-6.9, 0.95) {};
  \node[font=\sffamily\bfseries\scriptsize, text=pc] at (-6.9, 1.42) {Blender 5.2 LTS};
  \node[box=pc, fill=white, minimum width=2.3cm] (bla) at (-6.9, 0.8) {Extension Wysp\\\tiny bpy, timers, gpu};
  \node[box=pc, minimum width=2.6cm, minimum height=1.35cm] (ch) at (-6.9,-0.75) {};
  \node[font=\sffamily\bfseries\scriptsize, text=pc] at (-6.9,-0.28) {Chrome};
  \node[box=pc, fill=white, minimum width=2.3cm] (cha) at (-6.9,-0.9) {Extension MV3\\\tiny service worker, script};

  % Panel
  \node[box=pb, minimum width=2.6cm, minimum height=2.85cm] (panel) at (6.8, 0.12) {};
  \node[font=\sffamily\bfseries\scriptsize, text=pb] at (6.8, 1.3) {Panneau (React)};
  \node[box=pb, fill=white, minimum width=2.3cm] at (6.8, 0.75) {Tableau\\\tiny React Flow, ELK};
  \node[box=pb, fill=white, minimum width=2.3cm] at (6.8, 0.0) {Sous-titres, voix};
  \node[box=pb, fill=white, minimum width=2.3cm] at (6.8,-0.75) {Permissions\\\tiny exercices, leçons};

  % External
  \node[box=okgreen, minimum width=2.4cm] (prof) at (-1.9,-2.75) {Profils YAML\\\tiny apps/*/profile};
  \node[box=muted, minimum width=2.4cm] (prov) at (0.75,-2.75) {Fournisseur LLM\\\tiny appels d'outils};
  \node[box=muted, minimum width=2.4cm] (tts) at (3.65,-2.75) {Synthèse vocale\\\tiny clé côté serveur};

  % Arrows
  \draw[biarr] (bla.east) -- (gw.west |- bla.east);
  \draw[biarr] (cha.east) -- (gw.west |- cha.east);
  \node[font=\sffamily\tiny, text=muted, align=center] at (-5.05,-0.05) {JSON-RPC\\WebSocket};
  \draw[biarr] (ui.east |- 0,0.75) -- node[lab, above] {/ws/ui} (panel.west |- 0,0.75);
  \draw[arr] (prof.north) -- (ped.south);
  \draw[arr] (llm.south) -- (llm.south |- prov.north);
  \draw[arr] (ui.south) -- (tts.north);
  \node[lab] at (0.42,-1.78) {HTTPS};
  \node[lab] at (4.12,-1.78) {HTTPS};
\end{tikzpicture}
\caption{Vue d'ensemble. Seuls les adaptateurs (orange) et les profils (vert) dépendent de l'application.}
\end{figure}

\begin{xltabular}{\textwidth}{@{}L{3.1cm}L{2.2cm}Y L{1.2cm}@{}}
\toprule
\textbf{Composant} & \textbf{Langage} & \textbf{Rôle} & \textbf{Resp.} \\
\midrule
Contrat & Python (Pydantic) & Modèles, export JSON Schema, prédicats, empreinte d'état & \PA \\
Wysp Core & Python 3.13 & Passerelle, session, orchestrateur, LLM, tableau, narration, confiance, pédagogie, vérification & \PA \\
Kit de contrat, adaptateur factice & Python & Tests contre chaque adaptateur, application Kanban fictive & \PA \\
Panneau & TypeScript, React 19 & Tableau, sous-titres, voix, dialogues & \PB \\
Tableau de bord de démo & TypeScript, React & Application web cible « Haunted Shop Analytics » & \PB \\
SDK adaptateur & Python et TypeScript & Client WebSocket, registre d'actions, validation & \PC \\
Tuteur Blender & Python (bpy) & Extension Blender 5.2 et scène de démo (\PC), contenu du profil (\PB) & \PC{} \PB \\
Tuteur web & TypeScript & Extension Chrome MV3 (\PC), profil web (\PB) & \PC{} \PB \\
\bottomrule
\end{xltabular}

\subsection{Principes directeurs}

\begin{enumerate}[leftmargin=1.2em]
  \item \textbf{Le cœur ignore les applications.} Aucun import d'hôte (\code{bpy}, API du navigateur) ni vocabulaire d'application dans \file{core/}. Un script vérifie cette règle à chaque modification, en CI et via un hook Kiro.
  \item \textbf{Le LLM décide quoi dire et quoi montrer ; le cœur décide ce qui est vrai et ce qui est permis.} Le LLM n'a aucun outil d'exécution directe. Il propose ; la porte de permission décide ; l'adaptateur exécute ; le vérificateur contrôle.
  \item \textbf{Une seule source de vérité pour le contrat.} Modèles Pydantic, puis JSON Schema, puis types TypeScript générés. La CI échoue si les fichiers générés ne correspondent plus.
  \item \textbf{Les connaissances d'application sont des données.} Le profil (YAML) décrit compétences, concepts, leçons et exercices. Les prompts du cœur restent génériques.
  \item \textbf{Tout est testable sans l'application réelle} : adaptateur factice, LLM scénarisé, enregistrements d'échanges.
  \item \textbf{Local d'abord.} Écoute sur \code{127.0.0.1}, jeton d'appairage, liste blanche d'origines, clés côté serveur.
\end{enumerate}

\subsection{Le contrat d'adaptateur}

\subsubsection{Transport et session}

Les adaptateurs sont \textbf{clients} et le cœur est \textbf{serveur}. Le Python embarqué de Blender et le service worker d'une extension Chrome savent ouvrir une connexion WebSocket sortante, alors qu'héberger un serveur y serait difficile. Les messages sont du JSON-RPC 2.0 dans les deux sens sur la même connexion.

\begin{enumerate}[leftmargin=1.2em]
  \item L'adaptateur envoie \code{wysp.hello \{manifest, token\}}. Le cœur vérifie le jeton (fichier \file{~/.wysp/token}), la version majeure du contrat et le manifeste.
  \item L'adaptateur envoie \code{wysp.heartbeat} au moins toutes les 20\,s. Cette fréquence maintient aussi en vie le service worker Chrome (Chrome 116 et suivants). Le cœur déclare l'adaptateur déconnecté après 45\,s de silence.
  \item Le cœur appelle les méthodes \code{adapter.*}. L'adaptateur envoie des notifications \code{event.*}.
\end{enumerate}

{\small\setlength{\tabcolsep}{4pt}
\begin{xltabular}{\textwidth}{@{}L{3.7cm}L{3.0cm}L{2.0cm}L{0.9cm}Y@{}}
\toprule
\textbf{Méthode} & \textbf{Résultat} & \textbf{Capacité} & \textbf{Budget} & \textbf{Correspondance avec la note de concept} \\
\midrule
\code{adapter.observe} & Observation (graphe d'entités) & \code{observe} & 1\,s & \code{observe()} \\
\code{adapter.get\_context} & Contexte compact ($\leq$ 8\,Ko) & \code{context} & 0,5\,s & \code{get\_context()} \\
\code{adapter.get\_actions} & Catalogue d'actions & \code{actions} & 0,5\,s & \code{get\_actions()} \\
\code{adapter.execute} & Résultat, jeton d'annulation & \code{execute} & 5\,s & \code{execute(action)} \\
\code{adapter.undo} & Résultat & \code{undo} & 5\,s & ajouté : réversibilité \\
\code{adapter.highlight} & accusé & \code{highlight} & 0,5\,s & ajouté : montrer dans l'application \\
\code{adapter.verify} & Preuves & \code{verify.native} & 2\,s & \code{verify(goal)}, optionnel \\
\code{event.state\_changed} & notification & \code{events.state} & -- & événements \\
\code{event.selection\_changed}\newline \code{event.user\_action}\newline \code{event.navigation} & notifications & \code{events.*} & -- & événements \\
\bottomrule
\end{xltabular}}

\code{verify(goal)} et \code{evaluate(exercise)} de la note de concept sont surtout réalisés \textbf{par le cœur} : il évalue des prédicats sur une observation fraîche. Un adaptateur peut fournir \code{adapter.verify} quand une vérification native est plus fiable (par exemple un calcul de géométrie dans Blender). Ce choix évite de réécrire la logique de vérification dans chaque adaptateur.

\subsubsection{Le modèle d'observation : un graphe d'entités}

Une observation contient des \textbf{entités} identifiées par une référence stable de type URI, avec un \code{kind} déclaré dans le manifeste, un parent, des propriétés JSON et une confiance entre 0 et 1. Elle contient aussi des \textbf{liens} typés et un \textbf{focus} (sélection ou élément actif).

\begin{lstlisting}[language={}]
{"observation_id": "o-118", "state_version": 42, "scope": {"level": "focus"},
 "entities": [
  {"ref": "blender://object/Pumpkin", "kind": "object", "label": "Pumpkin",
   "props": {"type": "MESH", "modifiers_count": 2}},
  {"ref": "blender://object/Pumpkin/modifier/Subdivision", "kind": "modifier",
   "label": "Subdivision", "parent": "blender://object/Pumpkin",
   "props": {"type": "SUBSURF", "index": 0, "levels": 2, "show_viewport": true}},
  {"ref": "blender://object/Pumpkin/modifier/Solidify", "kind": "modifier",
   "label": "Solidify", "parent": "blender://object/Pumpkin",
   "props": {"type": "SOLIDIFY", "index": 1, "thickness": 0.02}}],
 "links": [], "focus": ["blender://object/Pumpkin"], "truncated": false}
\end{lstlisting}

La même structure décrit une page web : \code{web://…/region\textasciitilde{}filters/combobox\textasciitilde{}region} de type \code{control}, avec \code{role}, \code{name}, \code{value} et \code{states}. Le cœur ne connaît ni les objets 3D ni les champs de formulaire : il manipule des entités. Les limites sont 500 entités et 256\,Ko par observation, avec un indicateur \code{truncated} et des comptes par type.

\subsubsection{Actions, risques et post-conditions}

Chaque adaptateur publie un catalogue \textbf{fermé} d'actions. Une action déclare un identifiant, une description destinée au LLM, un schéma JSON de paramètres, un niveau de risque, des post-conditions et, si l'utilisateur peut la faire lui-même, un chemin d'interface (\code{guide.ui\_path}) utilisé en mode guidé. Aucune action ne peut évaluer du code arbitraire.

\begin{xltabular}{\textwidth}{@{}L{2.3cm}Y L{4.3cm}@{}}
\toprule
\textbf{Risque} & \textbf{Politique par défaut} & \textbf{Exemples} \\
\midrule
\code{read\_only} & Exécution automatique, journalisée & lecture de propriétés \\
\code{ui\_state} & Exécution automatique, journalisée & sélectionner un objet, changer d'onglet, faire défiler \\
\code{reversible} & Approbation : clic, « oui » vocal ou autorisation limitée à la leçon. Jeton d'annulation obligatoire & ajouter un modificateur, remplir un champ, choisir un filtre \\
\code{irreversible} & Clic explicite avec aperçu, jamais par la voix, interdit au niveau visuel & appliquer un modificateur, cliquer un bouton quelconque \\
\bottomrule
\end{xltabular}

\subsubsection{Langage de prédicats}

Les post-conditions, les critères de réussite des exercices et les objectifs de vérification utilisent un seul langage, évalué par le cœur : des comparaisons sur une expression JMESPath, combinées par \code{all}, \code{any} et \code{not}.

\begin{lstlisting}[language={}]
{"all": [
  {"path": "entities[?kind=='modifier' && parent=='${bindings.target}' && props.type=='BEVEL'] | length(@)",
   "op": "gte", "value": 1},
  {"path": "entities[?kind=='modifier' && parent=='${bindings.target}' && props.type=='BEVEL'].props.segments | [0]",
   "op": "eq", "value": 3}]}
\end{lstlisting}

Les paramètres ne sont substitués que dans des positions autorisées, avec échappement, pour empêcher une valeur d'injecter de la syntaxe JMESPath (propriété CON-P11). L'évaluateur renvoie la valeur réelle et la valeur attendue : ce sont les éléments de la \textbf{preuve} affichée à l'utilisateur.

\subsubsection{Empreinte d'état, erreurs et versions}

\begin{itemize}[leftmargin=1.2em]
  \item \textbf{Empreinte d'état} : SHA-256 d'un JSON canonique des entités (sans propriétés volatiles, flottants arrondis) et des liens triés. Elle définit « l'état est restauré » pour l'annulation et « rien n'a changé » pour les appels invalides.
  \item \textbf{Erreurs} : codes JSON-RPC standard et codes Wysp de \code{-32010} à \code{-32019} (non supporté, action inconnue, cible introuvable, précondition, échec, annulation impossible, délai, occupé, version incompatible, non autorisé).
  \item \textbf{Versions} : \code{contract\_version} commence à 1.0.0 au jalon J0. Ajout optionnel : version mineure. Changement cassant : version majeure, accord des trois personnes.
\end{itemize}

\subsubsection{Kit de tests de contrat}

Le kit (\code{wysp-contract --tutor <nom>}) joue le rôle du cœur : il démarre un serveur de test, lance l'adaptateur selon \file{apps/<tuteur>/contract.fixture.yaml}, puis vérifie la poignée de main, le manifeste, les invariants des observations, le catalogue, les appels sûrs (exécution, post-conditions, annulation, empreinte), les appels invalides, l'étiquetage d'origine des événements et les budgets de temps. Il enregistre les échanges. En mode \code{--offline}, il revalide ces enregistrements contre les schémas courants : chacun sait ainsi quel tuteur casserait un changement de contrat, sans lancer Blender ni Chrome.

\subsection{Wysp Core}

\subsubsection{Modules}

\begin{xltabular}{\textwidth}{@{}L{2.7cm}Y@{}}
\toprule
\textbf{Module} & \textbf{Responsabilité} \\
\midrule
\file{contract/} & Modèles Pydantic du contrat et du protocole du panneau, export des schémas, prédicats, empreinte d'état. Seul module importé par les SDK et le kit. \\
\file{gateway/} & Serveur JSON-RPC \code{/ws/adapter}, \code{AdapterProxy} typé avec budgets de temps, registre des adaptateurs, surveillance des battements, contrôle d'origine et de jeton. \\
\file{session/} & Session de tutorat (adaptateur actif, profil, tableau, narration, propositions), bus d'événements asyncio, mémoire de conversation avec résumé glissant, constructeur de contexte. \\
\file{orchestrator/} & Boucle de tour (\code{TurnRunner}), outils offerts au LLM, prompts génériques, annulation. \\
\file{llm/} & Interface \code{LLMProvider}, \code{ScriptedLLM} (scénarios rejouables), fournisseur compatible OpenAI (OpenAI, Groq, Ollama local…), Bedrock en option. \\
\file{board/} & Modèle visuel, opérations et correctifs versionnés, projection des entités, validateur de provenance, regroupement au-delà de 60 nœuds. \\
\file{narration/} & Segments, repères de synchronisation, curseur de lecture, interruption et reprise. \\
\file{trust/} & Politique de risque, porte de permission, autorisations de leçon, pile d'annulation, journal d'audit. \\
\file{verification/} & Observation fraîche, évaluation des prédicats, preuves. \\
\file{pedagogy/} & Profils, graphe de compétences, leçons, exercices, modèle de l'apprenant, recommandations. \\
\file{ui_api/}, \file{store/} & Protocole du panneau, proxy de synthèse vocale, page d'appairage ; persistance JSON/JSONL dans \file{~/.wysp/} avec masquage des secrets. \\
\bottomrule
\end{xltabular}

\subsubsection{La boucle de tour et les outils du LLM}

Un \textbf{tour} traite un message de l'utilisateur. Le cœur construit un contexte d'au plus 6\,000 jetons : prompt système générique, manifeste (types d'entités, catalogue d'actions), amorce du profil, contexte de l'adaptateur et observation du focus entourés de balises \code{<app\_data>} (contenu traité comme des données, jamais comme des instructions), état de l'apprenant, leçon en cours et conversation récente. Le LLM répond uniquement par des appels d'outils, au plus 8 tours d'outils par message.

\begin{xltabular}{\textwidth}{@{}L{3.0cm}Y L{2.3cm}@{}}
\toprule
\textbf{Outil} & \textbf{Effet} & \textbf{Capacité requise} \\
\midrule
\code{observe} & Résumé compact d'une observation & \code{observe} \\
\code{show\_entities} & Projection d'entités réelles sur le tableau (nœuds \pobs) & \code{observe} \\
\code{add\_concept}, \code{annotate} & Nœuds et annotations \pinf{} ou \pill, jamais \pobs & -- \\
\code{narrate} & Segments de narration avec repères (focus, surbrillance, étape) & -- \\
\code{highlight} & Mise en évidence dans l'application & \code{highlight} \\
\code{propose\_action} & Proposition à la porte de permission, mode \code{guide} ou \code{do} & \code{actions}, \code{execute} \\
\code{start\_exercise} & Démarrage d'un exercice du profil si ses liaisons se résolvent & \code{observe} \\
\code{lookup\_concept} & Entrée du glossaire du profil & -- \\
\code{ask\_clarification} & Question avec au plus 4 réponses proposées & -- \\
\bottomrule
\end{xltabular}

Les outils dont la capacité manque au manifeste ne sont pas proposés (propriété CORE-P10). Il n'existe aucun outil « exécuter » : c'est la garantie centrale du modèle de confiance.

\subsection{Modèle visuel et tableau}

Le tableau est un \textbf{modèle structuré}, pas une image : nœuds (avec \code{ref}, \code{kind}, provenance, confiance, propriétés à afficher), arêtes, étapes, annotations, focus et disposition. Le cœur publie des correctifs versionnés (\code{board.patch \{base\_version, version, ops\}}) ; le panneau les applique dans l'ordre et demande un instantané en cas de trou. La même fonction d'application existe en Python et en TypeScript, vérifiée par des fixtures communes (CORE-P3, PAN-P1).

\begin{itemize}[leftmargin=1.2em]
  \item \textbf{Projection.} Une entité devient un nœud d'identifiant \code{n:<ref>} ; les relations parent deviennent des arêtes \code{contains} ; les liens d'observation deviennent des arêtes typées. Le profil choisit les propriétés affichées par type (\code{board.props}). La projection est déterministe (CORE-P12), ce qui stabilise la disposition.
  \item \textbf{Validateur de provenance.} Tout nœud marqué \pobs{} dont la référence n'apparaît dans aucune observation de la session est rétrogradé en \pinf{} et journalisé (CORE-P1).
  \item \textbf{Dispositions.} Quatre suffisent pour les deux tuteurs et couvrent les exemples de la note de concept.
\end{itemize}

\begin{xltabular}{\textwidth}{@{}L{2.0cm}Y Y@{}}
\toprule
\textbf{Disposition} & \textbf{Blender} & \textbf{Web} \\
\midrule
\code{stack} & pile de modificateurs, ordre d'évaluation & ordre des filtres appliqués \\
\code{tree} & hiérarchie de la scène (collections, objets) & régions de la page, contrôles d'un formulaire \\
\code{graph} & objets, matériaux et références & dépendances entre filtres, KPI et tableau \\
\code{sequence} & étapes guidées d'une leçon & étapes de création d'un rapport \\
\bottomrule
\end{xltabular}

\begin{xltabular}{\textwidth}{@{}L{2.4cm}Y L{4.4cm}@{}}
\toprule
\textbf{Provenance} & \textbf{Style du nœud} & \textbf{Marqueur dans les sous-titres} \\
\midrule
\pobs & bordure pleine 2\,px, opacité complète, icône « œil » & aucun \\
\pinf & bordure pointillée, opacité 85\,\%, badge de confiance & « probablement » \\
\pill & fond hachuré, bordure grise, étiquette « exemple » & « par exemple » \\
\bottomrule
\end{xltabular}

\subsection{Voix et narration}

Nous retenons une chaîne en cascade plutôt qu'un modèle vocal de bout en bout : transcription dans le navigateur, raisonnement en texte dans le cœur, synthèse vocale par segments. Le texte reste ainsi la source de vérité : sous-titres, provenance et repères de synchronisation sont exacts, et la voix peut tomber en panne sans casser le tutorat.

\begin{figure}[h]
\centering
\begin{tikzpicture}[font=\sffamily\scriptsize, node distance=0.55cm,
  st/.style={draw=#1, fill=#1!8, rounded corners=2pt, align=center, minimum height=0.8cm, minimum width=1.9cm},
  arr/.style={-{Stealth[length=2mm]}, thick, color=surface}]
  \node[st=pb] (mic) {Micro\\\tiny appui pour parler};
  \node[st=pb, right=of mic] (stt) {Transcription\\\tiny Web Speech API};
  \node[st=pa, right=of stt] (core) {Cœur : tour\\\tiny LLM + outils};
  \node[st=pa, right=of core] (seg) {Segments\\\tiny $\leq$ 300 car. + repères};
  \node[st=pb, right=of seg] (tts) {Synthèse\\\tiny /api/tts ou navigateur};
  \node[st=pb, right=of tts] (spk) {Audio et\\sous-titres};
  \draw[arr] (mic) -- (stt); \draw[arr] (stt) -- (core); \draw[arr] (core) -- (seg); \draw[arr] (seg) -- (tts); \draw[arr] (tts) -- (spk);
  \draw[arr, color=danger, dashed] (spk.south) -- ++(0,-0.55) -| node[pos=0.25, below, text=danger] {interruption : parole détectée ou Échap $\rightarrow$ \code{user.interrupt}, arrêt en moins de 150\,ms} (mic.south);
\end{tikzpicture}
\caption{Chaîne vocale. En violet le panneau, en bleu le cœur.}
\end{figure}

\begin{itemize}[leftmargin=1.2em]
  \item \textbf{Appui pour parler} (barre d'espace ou bouton) par défaut : fiable en démonstration, pas de déclenchement intempestif.
  \item \textbf{Interruption} : pendant la narration, une détection d'énergie vocale (WebAudio, seuil calibré, annulation d'écho) ou la touche Échap arrête l'audio et envoie \code{user.interrupt}. Le cœur annule le tour et mémorise le dernier segment joué. « Continue » rejoue exactement les segments non entendus (CORE-P7, CORE-P8).
  \item \textbf{Clés} : la synthèse passe par \code{/api/tts} pour que la clé reste côté serveur ; la synthèse du navigateur sert de repli.
  \item \textbf{Latence visée} : premier son moins de 2,5\,s après la fin de la question, grâce au découpage en segments et à la prélecture du segment suivant.
\end{itemize}

\subsection{Pédagogie et vérification}

Le moteur pédagogique lit le \textbf{profil} de l'application : compétences avec prérequis, glossaire de concepts reliés à la documentation officielle, leçons en étapes (\code{explain}, \code{show}, \code{guide}, \code{do}, \code{check}) et exercices.

\begin{lstlisting}[language={}]
exercises:
  - id: bevel_segments
    prompt: {en: "Give the bevel 3 segments so the edge looks round."}
    skills: [bevel_modifier]
    bindings: {target: "$focus.object"}
    success: {all: [ ...predicates on the observation... ]}
    hints: [{en: "Look at the Modifier Properties tab (blue wrench)."},
            {en: "Find the Segments field in the Bevel panel."}]
    max_attempts: 3
\end{lstlisting}

\begin{itemize}[leftmargin=1.2em]
  \item \textbf{Exercices contextualisés} : un exercice n'est proposé que si ses liaisons se résolvent sur la scène ou la page réelle et si son critère n'est pas déjà satisfait (PED-P5, PED-P6).
  \item \textbf{Vérification} : à chaque changement d'origine utilisateur, le moteur évalue les prédicats sur une observation fraîche et enregistre une preuve (prédicat, valeur réelle, valeur attendue, identifiant d'observation, horodatage).
  \item \textbf{Trois niveaux de preuve} : \code{verified} (constaté dans l'application, poids 1,0), \code{answered} (réponse à un quiz, 0,6), \code{declared} (« c'est fait » non vérifiable, 0,3). Le panneau affiche toujours la différence. Les actions faites par Wysp ne comptent jamais (PED-P4).
  \item \textbf{Modèle de l'apprenant} : traçage bayésien des connaissances simplifié (glissement 0,1, chance 0,2, apprentissage 0,15), pondéré par le type de preuve et le nombre d'indices. Une compétence est acquise à 0,85. La profondeur des explications suit la maîtrise : bases sous 0,4, pratique guidée jusqu'à 0,85, défi au-delà.
\end{itemize}

\subsection{Confiance et sécurité}

\begin{figure}[h]
\centering
\begin{tikzpicture}[font=\sffamily\scriptsize,
  s/.style={draw=#1, fill=#1!8, rounded corners=2pt, align=center, minimum height=0.66cm, minimum width=1.75cm},
  arr/.style={-{Stealth[length=2mm]}, thick, color=surface},
  lbl/.style={font=\sffamily\tiny, text=muted, fill=white, inner sep=1pt}]
  \node[s=pa] (p) at (0,0) {proposée\\\tiny paramètres validés};
  \node[s=pa] (pol) at (2.4,0) {politique\\\tiny selon le risque};
  \node[s=okgreen] (auto) at (5.0,0.95) {approuvée\\\tiny auto ou autorisation};
  \node[s=pumpkin] (wait) at (5.0,-0.95) {en attente\\\tiny délai 60\,s};
  \node[s=muted] (deny) at (5.0,-2.35) {refusée\\\tiny ou expirée};
  \node[s=okgreen] (ex) at (7.7,0) {exécution\\\tiny adaptateur};
  \node[s=pa] (ver) at (10.1,0) {vérification\\\tiny post-conditions};
  \node[s=okgreen] (done) at (12.6,0.95) {réussie\\\tiny preuve enregistrée};
  \node[s=danger] (undo) at (12.6,-0.95) {échec\\\tiny annulation auto};
  \draw[arr] (p) -- (pol);
  \draw[arr] (pol) -- (auto);
  \draw[arr] (pol) -- (wait);
  \draw[arr] (auto) -- (ex);
  \draw[arr] (wait) -- node[lbl, pos=0.45, below right=-1pt] {clic ou « oui »} (ex);
  \draw[arr] (wait) -- node[lbl, right] {refus} (deny);
  \draw[arr] (ex) -- (ver);
  \draw[arr] (ver) -- (done);
  \draw[arr] (ver) -- (undo);
\end{tikzpicture}
\caption{Cycle de vie d'une proposition d'action.}
\end{figure}

\begin{itemize}[leftmargin=1.2em]
  \item \textbf{Provenance} : seule la projection d'observations crée des nœuds \pobs ; le validateur rattrape le reste. Au niveau d'intégration visuel, la confiance des entités est inférieure à 1 et affichée.
  \item \textbf{Porte de permission} : politique de la table des risques, autorisations limitées à une leçon pour les actions réversibles, clic obligatoire et aperçu pour les actions irréversibles, expiration à 60\,s. Une surcharge de politique par profil ne peut être que plus stricte (CORE-P4).
  \item \textbf{Vérifier ou annuler} : après exécution, nouvelle observation et évaluation des post-conditions ; en cas d'échec, annulation automatique et rapport avec preuve (CORE-P6). Chaque action Blender pousse aussi une étape d'annulation native, donc Ctrl+Z fonctionne.
  \item \textbf{Sécurité locale} : écoute sur \code{127.0.0.1} uniquement, contrôle de l'en-tête \code{Origin}, jeton d'appairage, aucune clé dans le navigateur, journaux masqués (CORE-P11).
  \item \textbf{Vie privée (web)} : extension active seulement sur les origines autorisées ; jamais de mot de passe, carte bancaire, code à usage unique, champ caché, cookie ou stockage dans les observations (WEB-P1, WEB-P6).
  \item \textbf{Injection de prompt} : le texte d'une page ou d'une scène est une donnée. Une page piégée peut au pire faire \emph{proposer} une action ; la porte et le catalogue fermé empêchent qu'elle s'exécute sans l'utilisateur.
  \item \textbf{Audit} : propositions, décisions (qui, par quel canal), exécutions, vérifications et annulations dans \file{~/.wysp/audit/*.jsonl} (CORE-P5).
\end{itemize}

\subsection{Adaptateur Blender (intégration native)}

L'adaptateur est une extension Blender 5.2 LTS (Python 3.13, \file{blender_manifest.toml}, permissions \code{network} et \code{files} justifiées, dépendances pures Python embarquées en wheels). Licence GPL-3.0-or-later pour ce dossier, comme l'exige l'écosystème des extensions Blender ; le reste du dépôt est sous MIT.

\begin{figure}[h]
\centering
\begin{tikzpicture}[font=\sffamily\scriptsize,
  t/.style={draw=#1, fill=#1!6, rounded corners=3pt, minimum width=5.6cm, minimum height=2.5cm},
  s/.style={draw=#1, fill=white, rounded corners=2pt, align=center, minimum height=0.55cm, minimum width=2.2cm},
  arr/.style={-{Stealth[length=2mm]}, thick, color=surface}]
  \node[t=muted] at (0,0) {}; \node[font=\sffamily\bfseries\scriptsize] at (0,1.0) {Fil réseau (websocket-client)};
  \node[t=pc] at (7.2,0) {}; \node[font=\sffamily\bfseries\scriptsize, text=pc] at (7.2,1.0) {Fil principal de Blender};
  \node[s=muted] (rx) at (-1.3,0.25) {réception\\\tiny requête JSON-RPC};
  \node[s=muted] (tx) at (-1.3,-0.65) {envoi\\\tiny réponse, événement};
  \node[s=muted] (hb) at (1.35,-0.2) {battement\\\tiny 15\,s, sans bpy};
  \node[s=pc, minimum width=2.5cm] (pump) at (5.75,0.25) {pompe \texttt{\tiny bpy.app.timers}\\\tiny 50\,ms, budget 20\,ms};
  \node[s=pc, minimum width=2.0cm] (h) at (8.75,0.25) {gestionnaire\\\tiny utilise bpy};
  \node[s=pc, minimum width=2.6cm] (dg) at (7.2,-0.65) {\texttt{\tiny depsgraph\_update\_post}\\\tiny anti-rebond 250\,ms};
  \draw[arr] (rx.east) -- node[above, font=\sffamily\tiny] {file d'entrée} (pump.west);
  \draw[arr] (pump) -- (h);
  \draw[arr] (dg.west) -- node[below, font=\sffamily\tiny] {file de sortie} (tx.east);
\end{tikzpicture}
\caption{Modèle de threads : \code{bpy} n'est jamais appelé hors du fil principal (BLD-P1).}
\end{figure}

\begin{itemize}[leftmargin=1.2em]
  \item \textbf{Observation} : entités \code{scene}, \code{collection}, \code{object}, \code{modifier}, \code{material}. Pour les modificateurs, extraction \emph{générique} des propriétés RNA modifiables de types simples (\code{mod.bl\_rna.properties}) : aucun code par type de modificateur, et tous les types de Blender 5.2 sont couverts (BLD-P3). Focus : objet actif et sélection.
  \item \textbf{Actions} : neuf actions, une par fichier. API de données plutôt qu'opérateurs (\code{obj.modifiers.new}, \code{obj.modifiers.move}) ; \code{temp\_override} quand un opérateur est inévitable.
  \item \textbf{Annulation} : enregistrement compensatoire (opération inverse) pour chaque action réversible, plus \code{bpy.ops.ed.undo\_push} pour que Ctrl+Z fonctionne aussi (BLD-P2).
  \item \textbf{Événements} : gestionnaire \code{depsgraph\_update\_post} avec anti-rebond, comparaison d'empreintes par entité, événements dérivés (« modificateur ajouté »), \code{bpy.msgbus} pour l'objet actif, réabonnement au chargement d'un fichier. Une fenêtre d'étiquetage autour de chaque action Wysp donne \code{origin: wysp} (BLD-P4).
  \item \textbf{Mise en évidence} : contour pulsé de la boîte englobante dans la vue 3D (module \code{gpu}), bascule de l'éditeur de propriétés sur l'onglet Modificateurs.
  \item \textbf{Mode sans interface} : \file{demo/run_adapter.py} lance l'adaptateur avec \code{blender -b} et une pompe manuelle, pour le kit de contrat et la CI.
\end{itemize}

{\small
\begin{xltabular}{\textwidth}{@{}L{5.0cm}L{1.9cm}Y@{}}
\toprule
\textbf{Action} & \textbf{Risque} & \textbf{Chemin d'interface (mode guidé)} \\
\midrule
\code{select\_object}, \code{show\_properties\_tab} & ui\_state & clic dans l'Outliner ; icône d'onglet de l'éditeur de propriétés \\
\code{add\_modifier} & reversible & Propriétés › Modificateurs (clé) › Add Modifier › catégorie › type \\
\code{set\_modifier\_property} & reversible & panneau du modificateur › champ \\
\code{remove\_modifier}, \code{move\_modifier} & reversible & menu de l'en-tête du modificateur \\
\code{toggle\_modifier\_visibility} & reversible & icône écran de l'en-tête \\
\code{assign\_material} & reversible & Propriétés › Matériau › emplacement \\
\code{apply\_modifier} & irreversible & menu › Apply (Ctrl+A), avec aperçu du nombre de sommets \\
\bottomrule
\end{xltabular}}

\subsection{Adaptateur web (intégration par l'interface)}

\begin{figure}[h]
\centering
\begin{tikzpicture}[font=\sffamily\scriptsize,
  s/.style={draw=#1, fill=#1!8, rounded corners=2pt, align=center, minimum height=0.9cm, minimum width=2.6cm},
  arr/.style={{Stealth[length=2mm]}-{Stealth[length=2mm]}, thick, color=surface}]
  \node[s=pa] (core) at (0,0) {Wysp Core\\\tiny /ws/adapter};
  \node[s=pc] (sw) at (4.2,0) {Service worker\\\tiny client SDK, battement 20\,s};
  \node[s=pc] (cs) at (8.4,0) {Script de contenu\\\tiny instantané, actions, overlay};
  \node[s=muted] (dom) at (12.4,0) {Page autorisée\\\tiny DOM, accessibilité};
  \draw[arr] (core) -- node[above, font=\sffamily\tiny] {JSON-RPC} (sw);
  \draw[arr] (sw) -- node[above, font=\sffamily\tiny] {runtime.Port} (cs);
  \draw[arr] (cs) -- (dom);
\end{tikzpicture}
\caption{Extension Chrome MV3 : le service worker tient la connexion, le script de contenu fait tout le travail sur le DOM.}
\end{figure}

\begin{itemize}[leftmargin=1.2em]
  \item \textbf{Instantané orienté accessibilité} : rôles ARIA et HTML, noms et descriptions accessibles calculés avec \code{dom-accessibility-api}. Types d'entités : \code{page}, \code{region}, \code{heading}, \code{form}, \code{control}, \code{dialog}, \code{table}, \code{list}, \code{chart}, \code{status}. Options des listes (20 au plus), en-têtes et 10 premières lignes des tableaux, libellé des graphiques.
  \item \textbf{Localisateurs et références} : rôle + nom accessible dans la région nommée la plus proche, puis index, puis \code{data-testid}, puis chemin CSS. La référence encode le localisateur, donc elle reste stable tant que l'élément garde son rôle, son nom et sa région (WEB-P2, WEB-P3). Une cible ambiguë renvoie \code{TARGET\_NOT\_FOUND} avec des candidats.
  \item \textbf{Actions} : \code{scroll\_into\_view}, \code{focus\_element}, \code{follow\_link} (ui\_state, même origine seulement), \code{fill\_field}, \code{select\_option}, \code{set\_checked} (réversibles, l'annulation restaure la valeur précédente), \code{click\_element} (irréversible, aperçu de la cible). Les valeurs sont posées avec le \emph{setter} natif suivi des événements \code{input} et \code{change}, ce qui fonctionne avec React.
  \item \textbf{Origine des événements} : les événements synthétiques de Wysp ont \code{isTrusted === false}, ceux de l'utilisateur \code{true}. Les mutations du DOM sont attribuées à Wysp pendant la fenêtre d'exécution d'une action (WEB-P5).
  \item \textbf{Overlay} : racine Shadow DOM fermée, \code{pointer-events: none}, ignorée par l'instantané et l'observateur de mutations (WEB-P4).
  \item \textbf{Cible de démonstration} : « Haunted Shop Analytics » (\PB), tableau de bord accessible avec KPI, graphique, filtres (dates, région, catégories), tableau des commandes, dialogue « Create report » et liste des rapports, données déterministes et bouton de réinitialisation.
\end{itemize}

\subsection{Pile technique et structure du dépôt}

\begin{xltabular}{\textwidth}{@{}L{3.0cm}Y@{}}
\toprule
\textbf{Domaine} & \textbf{Choix} \\
\midrule
Cœur & Python 3.13, uv (espace de travail), FastAPI + uvicorn, Pydantic v2, jmespath, PyYAML, httpx ; pytest, Hypothesis, ruff \\
Panneau & React 19, Vite, \code{@xyflow/react} (React Flow 12), elkjs, Zustand ; Vitest, Testing Library, fast-check ; eslint, prettier \\
Blender & Blender 5.2 LTS, extension, websocket-client en wheel ; \code{bpy} de PyPI et fake-bpy-module pour les tests et l'éditeur \\
Navigateur & Chrome 116+, Manifest V3, dom-accessibility-api, Vite ; Playwright \\
Services & LLM avec appels d'outils (fournisseur compatible OpenAI par défaut, Bedrock en option), Web Speech API pour la transcription, synthèse vocale par proxy \\
CI & GitHub Actions : tests Python et TypeScript, export des schémas sans différence, contrôle de frontière, kit de contrat sur l'adaptateur factice \\
\bottomrule
\end{xltabular}

\begin{lstlisting}[language={}]
.kiro/            steering, specs, hooks, agents (commun, versionne)
core/             Wysp Core, paquet wysp_core                          (A)
schemas/          JSON Schemas generes                                 (A)
contract-tests/   kit de contrat, adaptateur factice Kanban, fixtures  (A)
sdk/python/       SDK adaptateur Python                                (C)
sdk/ts/           SDK TypeScript + types generes                       (C, types : A)
panel/            panneau React                                        (B)
apps/blender-tutor/   adapter/ (extension), profile/, demo/            (C)
apps/web-tutor/       extension/ (C), demo-dashboard/ (B), profile/ (B)
docs/             ce document, guide d'adaptateur, journal Kiro
scripts/kiro/     hook.py : controles appeles par les hooks Kiro et la CI
\end{lstlisting}

\subsection{Décisions d'architecture}

\begin{xltabular}{\textwidth}{@{}L{3.2cm}L{3.6cm}Y@{}}
\toprule
\textbf{Décision} & \textbf{Alternatives écartées} & \textbf{Raison} \\
\midrule
JSON-RPC 2.0 sur WebSocket, adaptateurs clients & HTTP REST ; gRPC ; MCP comme protocole d'adaptateur & Bidirectionnel (événements), simple dans Blender et dans un service worker, aucun serveur à héberger dans l'hôte. MCP reste utile côté développement (pont en lecture seule), mais ne porte ni événements ni permissions. \\
Pydantic source unique, JSON Schema, types TS générés & Écrire les schémas à la main ; TypeScript source & Une seule définition, validation à l'exécution des deux côtés, contrôle de dérive en CI. \\
Graphe d'entités générique & Un modèle par application & Le cœur, le tableau et les prédicats restent indépendants de l'application. \\
Vérification par prédicats évalués dans le cœur & \code{verify()} écrit dans chaque adaptateur & Moins de code par adaptateur, preuves uniformes ; \code{adapter.verify} reste possible. \\
Nœuds « observés » créés par projection & Laisser le LLM écrire le tableau & Garantit la règle de confiance de la note de concept. \\
Voix en cascade, appui pour parler & Modèle vocal temps réel de bout en bout & Sous-titres et provenance exacts, interruption maîtrisée, coût et dépendance moindres, repli possible. \\
Panneau web dans une fenêtre Chrome & Panneau natif dans Blender ; Electron & Une seule interface pour les deux tuteurs, c'est la preuve visuelle du cœur commun. \\
Catalogue d'actions fermé & Laisser le LLM écrire du \code{bpy} ou du JavaScript & Sécurité, vérifiabilité, annulation possible. \\
Stockage local en fichiers & Base de données & Suffisant pour un prototype, transparent pour l'utilisateur. \\
\bottomrule
\end{xltabular}

